\documentclass[11pt,twoside]{amsart}
\usepackage{amssymb, amsmath, enumerate, libertine, hyperref,xcolor,booktabs,longtable,microtype,multirow}
\usepackage[normalem]{ulem}
\usepackage{fullpage}
\usepackage[T1]{fontenc}
\renewcommand{\labelitemi}{$\cdot$}
\usepackage[normalem]{ulem}
\definecolor{dark-red}{rgb}{0.4,0.15,0.15}
%   \definecolor{dark-blue}{rgb}{0.15,0.15,0.4}
%   \definecolor{medium-blue}{rgb}{0,0,0.5}
\setcounter{secnumdepth}{2}
\setcounter{tocdepth}{1}
\hypersetup{
    colorlinks, linkcolor=dark-red,
    citecolor=dark-red, urlcolor=dark-red
}

\title[Math 111:  Calculus]{Math 111: Calculus F03 \& F04\\ Course Information \& Syllabus}
\author[Math 111:  Calculus]{Fall 2026}

\begin{document}
\maketitle

%\vspace{-5mm}
\thispagestyle{empty}

\vspace{-.5cm}

\begin{center}
\fbox{
\begin{minipage}{4.5in}
\begin{tabular}{rl}
Place: &F03 --- Eliot 207\\
&F04 --- Chem 301\\
Time: &F03 --- 10--10:50\textsc{a.m.}\\
&F04 --- 12--12:50\textsc{p.m.}\\
Instructor: &Kyle Ormsby (\href{mailto:ormsbyk@reed.edu}{ormsbyk@reed.edu})\\
Office Hours: &MW 1:00--2:30\textsc{p.m.}, Library 306\\
&or \href{https://calendar.app.google/xUQmAFHKmqt44Q7R8}{book a Friday appointment}\\
Course Assistant: &Ezra Rosenfield\\
Problem Sessions: &Th 4:40--6:40\\textsc{p.m.}, Library 389\\
Textbook: &\href{https://activecalculus.org/acs2e/}{Active Calculus} (2nd ed.)\\
Website: &\href{https://e-infinity.space/111/}{e-infinity.space/111/}\\
Zulip: &\href{https://math111fall26.zulipchat.com/}{math111fall26.zulipchat.com}
\end{tabular}
\end{minipage}
}
\end{center}

\smallskip

\subsection*{Course description}
Calculus --- the mathematical study of change --- is one of humanity's most fruitful inventions and is fundamental to the study of physical sciences, computer algorithms, and technology. This course is an introduction to differential and integral single variable calculus, with a focus on \emph{what the ideas mean} and on reasoning carefully about them, rather than on accumulating computational techniques. We will explore fundamental mathematics and wide-ranging applications, and we will invest significant time in developing careful reasoning and communication. There are no formal proofs in this course.

\subsection*{Course goals}
Below are the five things I want you to be able to do at the end of this course. They are not equally large, and I have not tried to make them so. Everything we do is in service of one or more of them.
\begin{itemize}
\item \textbf{G1 --- Rate.} Given a quantity that varies, you can say what its rate of change measures, in what units, and what a particular value of that rate tells you about the situation --- whether the quantity is given by a formula, a graph, a table, or only words.
\item \textbf{G2 --- Accumulation.} You can recognize when a question is asking what has piled up rather than how fast something is going, explain what a definite integral adds up and in what units, and set up the integral that answers the question.
\item \textbf{G3 --- The connection.} You can explain why rate and accumulation are inverse problems, state both Fundamental Theorems in your own words, and say what each one lets you do that the other doesn't.
\item \textbf{G4 --- Approximation and error.} For any of these quantities you can produce an over-estimate and an under-estimate, bound your error, and say what you would do to bring the error below a stated tolerance.
\item \textbf{G5 --- Reasoning and communication.} You can decide whether a theorem's hypotheses are met and what follows when they aren't; produce examples and counterexamples to order; find the first wrong step in a flawed argument and repair it; and write an explanation another student in the course could read and follow.
\end{itemize}
Underneath all five sits \textbf{fluency}: differentiating polynomial, rational, exponential, logarithmic, and trigonometric functions; the power, product, quotient, and chain rules; antidifferentiation and substitution; and evaluating definite integrals using the Fundamental Theorem. Fluency is not an end in itself --- you cannot reason about an object you cannot manipulate --- but you will need it, and our exams will ask for it.

\subsection*{Distribution requirements}
This course can be used towards your Group III, ``Natural, Mathematical, and Psychological Science,'' requirement.  It accomplishes the following goals for the group:
\begin{itemize}
\item Use and evaluate quantitative data or modeling, or use logical/mathematical reasoning to evaluate, test, or prove statements.
\item Given a problem or question, formulate a hypothesis or conjecture, and design an experiment, collect data or use mathematical reasoning to test or validate it.
\end{itemize}
This course \textbf{does not} satisfy the ``primary data collection and analysis'' requirement.

\subsection*{Text}
We will use \emph{Active Calculus} (2nd ed.) as a reference for the course.  The book is freely accessible  online at \href{https://activecalculus.org/acs2e/}{activecalculus.org/acs2e/}.  Each lecture will be paired with a \textbf{required reading} that you need to do before class.  Lectures and readings will be complementary, and both will be necessary to fully appreciate and learn the course content. Note that we will not read the book straight through: we take up accumulation before antidifferentiation, and we skip several sections entirely. The course website has the authoritative reading schedule.

You will also need copies of the \emph{Active Calculus Activities Workbooks}. You can either write in print versions of the workbooks (available for purchase at the bookstore) or use a tablet computer to write on the PDF versions of the workbooks:
\begin{itemize}
\item \href{https://activecalculus.org/wp-content/uploads/2024/08/acs-activity-workbook-14-2024.pdf}{Ch.~1--4 Workbook} [pdf]
\item \href{https://activecalculus.org/wp-content/uploads/2024/08/acs-activity-workbook-58-2024.pdf}{Ch.~5--8 Workbook} [pdf]
\end{itemize}

I also \textbf{strongly recommend} doing the auto-assessed WeBWorK exercises at the end of each section of our textbook. (You'll need to use the HTML version of the book to access these.) These sets stay open all semester with unlimited attempts, and they are strictly for practice --- which is important! --- but will not count towards your final grade.

Before each exam I will also post a \textbf{mixed practice set}: a shuffled collection of problems drawn from everything we have done so far, with no labels telling you which section a problem comes from. These will feel harder than the section-by-section sets, and that is the point. On an exam, nobody tells you which tool to reach for; deciding \emph{which} technique a problem calls for is a separate skill from executing it, and the only way to practice it is on problems that arrive unsorted.

Finally, you can find video lectures on each section of the book on GVSUmath's YouTube channel:
\begin{itemize}
\item \href{https://www.youtube.com/playlist?list=PL9bIjQJDwfGuXQHuS5Jkmum_CFILoCZX-}{Chapters 1--4},
\item \href{https://www.youtube.com/playlist?list=PL9bIjQJDwfGtewW75Nw7PnGNSkfqwAm3v}{Chapters 5--6}.
\end{itemize}


\subsection*{Participation}
I will deliver interactive in-person lectures with frequent breaks for collaborative problem-solving. Your required reading will lay the groundwork for productively participating in these meetings. To demonstrate that you have completed the reading, and also allow me to adapt our course meetings to student needs, you are required to submit daily reading reflections via \href{https://docs.google.com/forms/d/e/1FAIpQLSdxdgkPt8zIwbsTJmZi3FsvcsQxY0sLmPuc_-qDOIrcYaZ_jg/viewform}{this Google Form}.

Submitting reading reflections, attending class, and participating in collaborative problem-solving will form the \textbf{engagement} part of your course grade. If you miss a class, it is your responsibility to catch up with the material via the course notes, textbook, office hours, and discussions with peers.

All students are encouraged --- but not required --- to attend office hours and problem sessions (see below).

\subsection*{Exams}\label{exams}
We will have \textbf{three in-class exams}, on \textbf{Friday, September 25}, \textbf{Friday, October 16}, and \textbf{Friday, November 20}. The meeting before each exam is a \textbf{review day}. Review days are for synthesis, not drill: we will spend them putting the unit's ideas next to each other and asking how they fit together. Come with questions, and in particular come with the problems from the mixed practice set. I will post a practice exam a week in advance.

Our exams are \textbf{cumulative and nested}. The second exam includes material from the first part of the course; the third includes everything before it; the final includes everything. This is deliberate. Calculus is not a sequence of independent techniques to be learned and discarded, and revisiting earlier material repeatedly --- rather than once, intensively, the night before --- is the single best-supported thing you can do to actually retain it.

Exams are closed-book and device-free, so it is important that you develop unaided problem-solving skills through the practice problems and written assignments.

\subsection*{Written assignments}\label{hw}
Roughly every other week you will complete a written assignment. There are five of them, due via Gradescope by 10:00\textsc{p.m.}~on \textbf{September 18}, \textbf{October 9}, \textbf{October 30}, \textbf{November 13}, and \textbf{December 4}. Each consists of three or four questions requiring synthesis and deeper thought: interpreting a derivative or integral in context, reasoning from a graph or a table with no formula available, deciding whether a theorem applies, constructing an example to order, or finding the first wrong step in a flawed argument and repairing it.

In your solutions, you are expected to communicate mathematics clearly in sentences, paragraphs, and notation. Excellent solutions take many forms, but they all have the following characteristics:
\begin{itemize}
\item they are written as explanations for other students in the course; in particular, they fully explain all of their reasoning and do not assume that the reader will fill in details;
\item when graphical reasoning is called for, they include large, carefully drawn and labeled diagrams;
\item they are neatly written or typeset;\footnote{Interested students are
encouraged to prepare solutions in the \LaTeX~document preparation
system.  Nearly all of the \texttt{.pdf} files on the course website are produced by \LaTeX; you can find their associated source files by changing the \texttt{.pdf} suffix to \texttt{.tex} in the file's URL.} and
\item they use complete sentences, even when formulas or symbols are involved.
\end{itemize}

Correctly answering these questions is how you \textbf{synthesize and communicate} your learning in calculus. You are permitted and encouraged to work with your peers on these problems.  You must cite those with whom you worked, and you must write up solutions independently.  \textbf{Duplicated solutions will not be accepted and constitute a violation of the Honor Principle.} I recommend starting early so that you can take advantage of office hours and problem sessions.

\subsection*{Oral exam}
Late in the semester each of you will have a short individual conversation with me about calculus --- roughly fifteen minutes, scheduled outside of class time (dates TBD). A week beforehand I will circulate six questions. At the start of your slot you will roll a die to select three of them, and we will talk through those, with follow-up questions from me. There is nothing to prepare beyond preparing to talk about the course and gaining confidence with the assigned problems.

This is the part of the course where I find out what you actually understand, and it is worth being clear about why it exists. Written work tells me what you can produce; a conversation tells me what you can explain. The oral does not carry a percentage. Instead it adjusts your final letter grade by at most one step in either direction (see below).

\subsection*{Feedback}
You will receive timely feedback on your exams and written assignments via Gradescope, either from me or the course grader (a mathematics undergraduate).

Each problem is assessed \textbf{holistically} --- as a whole piece of mathematical thinking, not as a checklist of steps --- on the following scale:
\begin{itemize}
\item \textbf{5}: perfect.
\item \textbf{4}: minor mistakes.
\item \textbf{3}: major mistake, right idea.
\item \textbf{2}: wrong, but contains a significant idea.
\item \textbf{1}: wrong, but contains a relevant idea.
\item \textbf{0}: none of the above.
\end{itemize}
Written assignment problems additionally receive up to \textbf{two points for the quality of the writing}, assessed against the characteristics listed above. Communication is a course goal, not a courtesy, and it is graded as one.

\subsection*{Deadlines, extensions, and revisions}
Learning takes time and consistency. The deadlines and flow of work in this course are designed to facilitate your learning. I acknowledge, though, that sometimes your time won't cooperate with the course schedule. Each student therefore has \textbf{five no-questions-asked 48-hour extensions} on written assignments, requested through \href{https://docs.google.com/forms/d/e/1FAIpQLSfIfgx44SFtUzefTNr-80yz35MVJaUqQttt96xVTU5mCQCOgg/viewform}{this form}. You do not need to explain yourself; you just need to ask before the deadline. Extensions cannot be applied to reading reflections, exams, or the oral exam.

Each written assignment also comes with one \textbf{revision window}. After an assignment is returned, you may rewrite any problem on it and resubmit; revised work is regraded on the same scale and the revised mark replaces the original. Revisions are due two weeks after the assignment is returned --- a fixed date, announced when I hand the work back, not a rolling deadline.

Revision is where most of the learning in a writing-heavy course actually happens. Please use it, and please use it on the problems you got \emph{wrong} rather than the ones you nearly got right.

\subsection*{Final exam}
We will have a comprehensive final exam as scheduled by the registrar during the December 14--17 exam period. It covers the entire course.

The final also functions as a safety net: \textbf{if your score on the final exceeds your lowest in-class exam score, it replaces that score}. A bad day in September, or a unit that only clicked in November, need not follow you to the end of the term.

\subsection*{What each kind of work is for}
Your work is assessed in four ways, and each one is asking a different question.  The
table below says which goals each is aimed at and why that work happens there.  It is worth
reading before the grade weights below, because the weights are downstream of these
decisions rather than the other way around.

{\small
\begin{longtable}{@{}p{1.5in}p{1.35in}p{3.05in}@{}}
\toprule
\textbf{Channel} & \textbf{Probes} & \textbf{Why there}\\
\midrule
\endfirsthead
\multicolumn{3}{@{}l}{\emph{What each kind of work is for, continued}}\\[2pt]
\toprule
\textbf{Channel} & \textbf{Probes} & \textbf{Why there}\\
\midrule
\endhead
\bottomrule
\endlastfoot
Cumulative in-class exams, device-free
& G1--G4, and fluency throughout
& Holistic judgment on whole pieces of work.  Because each exam includes earlier
  material, fluency is spaced and re-tested without a separate quiz track.\\
\addlinespace
Written assignments, low weight and revisable
& G5 chiefly, G4 alongside it
& Where argument and prose live, and where having time to think is the point.\\
\addlinespace
The oral exam, which adjusts the letter grade
& All five --- above all, whether you can \emph{say} it
& Talking through a question is the only place either of us finds out whether the
  understanding is real.\\
\addlinespace
Engagement: reading reflections, class participation
& Preparation
& Built on what you actually submit rather than an impression of participation, and it
  tells me what needs attention that day.\\
\end{longtable}}

\noindent Fluency has no row of its own.  That is deliberate: it is the condition for
everything else rather than a fifth thing to be graded, and the exams assess it by
requiring you to use it.

\subsection*{Grades}\label{grading}
Your final grade in this course is based on the following:
\begin{center}
\begin{tabular}{@{}lr@{}}
\toprule
Three in-class exams (15\% each) & 45\%\\
Final exam & 25\%\\
Written assignments & 20\%\\
Engagement (reflections, attendance, participation) & 10\%\\
\bottomrule
\end{tabular}
\end{center}

\medskip

The oral exam does not carry a percentage. It adjusts the letter grade produced by the table above by at most one step in either direction (\emph{e.g.}, B to B+, or B to B--). The grades of D and F do not admit $+$ or $-$ adjustments at Reed, and thus these modifications do not apply to those grades.

These weights determine a \emph{provisional} grade. To produce your final grade, I take into account the whole record of your work in the course. The question I am trying to answer is whether you understand calculus, not whether you have accumulated points. In practice, this discretion runs in one direction --- I reserve the right to adjust a grade upward when the number does not reflect what I have seen from you, and I will not adjust one downward below what the weights produce.

If this sounds less precise than a formula, that is because it is. Grades at Reed are a professional judgment made by your instructor, and I would rather make that judgment honestly than hide it behind arithmetic that only looks objective.

\subsection*{Joint expectations}
As members of a communal learning environment, we should all expect consideration, fairness, patience, and curiosity from each other.  Our aim is to all learn more through cooperation and genuine listening and sharing, not to compete or show off.  I expect diligence and academic and intellectual honesty from each of you.  You should expect that I will do my best to focus the course on interesting, pertinent topics, and that I will provide feedback and guidance that will help you excel as a student.

\subsection*{Help}
There are a number of resources you can access for help with this course's content.

My drop-in office hours are Mondays and Wednesdays, 1:00--2:30\textsc{p.m.}, in Library 306.  No appointment is needed; come with a question, or just come to work in the room.  Office hours are an opportunity to clarify difficult material and also to delve deeper into topics that interest you, and much of what you want to ask is already on a classmate's mind---you will often get more out of a shared hour than a private one.

If you would rather meet one-on-one, you can \href{https://calendar.app.google/xUQmAFHKmqt44Q7R8}{book a 20-minute appointment} on Fridays between 1:00 and 3:00\textsc{p.m.}, also in Library 306.  Book at least 24 hours in advance, up to two weeks out.

\textbf{If none of these times fit your schedule, email me at \href{mailto:ormsbyk@reed.edu}{ormsbyk@reed.edu} and we will find a time that does.}  This is a standing offer, not a last resort: work, athletics, care responsibilities, and courses that collide with my hours are all ordinary reasons to ask.  Please reach out if there are any other barriers preventing you from making use of these opportunities.

Our course assistant, Ezra Rosenfield, will run a two-hour problem session each week, Thursdays 4:40--6:40\\textsc{p.m.} in Library 389.  We will share problem sessions with the other sections of Math 111.  The problem sessions will provide a structured, facilitated environment in which you can collaborate on homework.  All students are encouraged to attend.

The Mathematics \& Statistics Department also hosts drop-in tutoring Sunday through Thursday, 6:00--8:00\\textsc{p.m.} in Library 204. Upperclass tutors will be available to clarify concepts and help you with homework problems.

Finally, every Reed student is entitled to one hour of free individual tutoring per week.  Use the tutoring app in IRIS to arrange to work with a student tutor.

I also maintain the \href{https://e-infinity.space/math-th/}{\emph{Math Trailhead}}, a
self-paced review of the algebra, functions, and trigonometry this course
assumes. It is organized by topic rather than by course, so you can start
wherever you feel shaky; for Math 111 I suggest beginning with
\href{https://e-infinity.space/math-th/functions.html}{Functions} and
\href{https://e-infinity.space/math-th/algebra.html}{Algebra}. Every exercise is
a live WeBWorK problem with immediate feedback and a fresh version whenever you
want more practice, and nothing you enter is recorded.

\subsection*{Zulip}
Our sections of Math 111 have a Zulip workspace.  Use the Zulip workspace to ask questions (of me or the class), collaborate on problems, and share resources. The Zulip workspace is an extension of our classroom and the above joint expectations extend to this setting.

You will receive an email invitation to join \href{https://math111fall26.zulipchat.com/}{our Zulip workspace} during the first week of classes. Please use channels and threads to keep conversations organized.

\subsection*{Technology}
The use of electronic devices (cell phones, computers, tablets, calculators, \emph{etc}.) is prohibited in the classroom without prior authorization from the instructor.  That said, legitimate uses of technology (\emph{e.g.}, note-taking or filling in the PDF versions of the workbooks) will be accommodated ---
just talk to me first. There will also be days when I ask you to bring a laptop or tablet, because we are using a graphing utility as an instrument; I will tell you in advance.

Exams and the oral exam are device-free.

\subsection*{The Internet and ``AI''}
You are welcome to use Internet resources to supplement content we cover in this
course, with the exception of solutions to assigned problems. \textbf{Copying
solutions from the Internet is an Honor Principle violation and will result in an
academic misconduct report.}

For large language models (``AI'' chatbots), what is permitted depends on which
work you are doing.

\textbf{On practice work} --- WeBWorK exercises, textbook problems, and the mixed
practice sets --- you may use a chatbot to have a concept explained a different
way, to generate extra problems, or to ask whether a final answer you have
already produced is correct. If it tells you that you are wrong, that is where
the conversation ends: work out \emph{why} yourself, or bring it to office hours,
the problem session, or Zulip. A model may tell you \emph{that} an answer is
wrong; it may not tell you what the right one is, or where your reasoning broke
down.

\textbf{On the five written assignments}, do not use a chatbot at all. These
problems ask you to construct examples, reason from graphs and tables, decide
whether a theorem applies, and repair flawed arguments --- there is usually no
``final answer'' to check, and the reasoning \emph{is} the submission. Asking a
model to solve, draft, outline, or critique one of these problems is asking it to
do the assignment.

If you used a chatbot anywhere in your work for this course, say so at the end of
your submission and include your prompts. Undisclosed use is an Honor Principle
violation.

Finally, \textbf{you must be able to explain anything you submit.} I may ask you
to walk me through a written assignment problem in person, without advance
notice. This is a short conversation, not a second oral exam, and it is graded
only in the sense that work you cannot explain receives no credit.

\subsection*{Academic accommodations}
If you have a documented disability requiring academic accommodation, please have  Disability \& Accessibility Resources (DAR)  provide a letter during the first week of classes.  I will then contact you to schedule a meeting during which we can discuss your accommodations.  If you believe you have an undocumented disability and that accommodations would ensure equal access to your Reed education, I would be happy to help you contact DAR.

\subsection*{Acknowledgments}
The structure and emphasis of this course benefits greatly from the conversations I've had with colleagues at Reed, especially Marcus Robinson, Jamie Pommersheim, and participants in Center for Teaching and Learning workshops. Some of our in-class activities are adapted, with thanks, from Peter Keep's \href{https://www.discovercalculus.com/}{\emph{Discover Calculus}}, shared under a \href{https://creativecommons.org/licenses/by/4.0/}{CC BY 4.0} license.

\subsection*{AI use disclosure}
I will use an LLM to help manage the course website and digital resources. All lectures, exams, and mathematical content will be prepared by me without the use of an LLM.

\bigskip \bigskip

\begin{center}
Remember: \emph{Math is hard, but we'll get through this together!}
\end{center}


\end{document}
